\section{Contexte sectoriel : Le conseil à l'ère de l'intelligence artificielle}
Le secteur du conseil en stratégie et en management traverse une période de mutation technologique sans précédent. Historiquement définies comme des \textit{Knowledge-Intensive Firms} (KIFs), les entreprises de conseil tirent leur valeur ajoutée de la capacité de leurs collaborateurs à collecter, analyser et synthétiser des informations complexes pour éclairer la prise de décision des dirigeants. 

L'émergence et la démocratisation fulgurante des outils d'intelligence artificielle générative depuis la fin de l'année 2022 sont venues bousculer ce modèle économique et opérationnel. Des tâches autrefois dévolues aux consultants juniors, telles que la revue documentaire, la synthèse de rapports sectoriels ou la première formulation de recommandations, peuvent désormais être automatisées ou fortement accélérées par des modèles de langage de grande taille (LLM). 

Cette transformation technologique offre des opportunités majeures en termes d'efficience opérationnelle et de gain de productivité pour les cabinets. Toutefois, elle soulève des questionnements managériaux profonds. D'une part, l'automatisation des tâches d'analyse de bas niveau risque d'altérer le processus traditionnel d'apprentissage et de développement des compétences des collaborateurs en début de carrière. D'autre part, la récurrence des biais algorithmiques et des hallucinations informationnelles pose un risque de réputation et de responsabilité juridique majeur pour des structures dont la "confiance" constitue le principal actif immatériel.

\section{Présentation du cadre de l'étude : Wavestone Luxembourg}
C'est au cœur de cette dynamique que s'inscrit ce travail de recherche, mené dans le cadre d'un stage de fin d'études au sein du cabinet Wavestone Luxembourg. Acteur majeur du conseil en transformation digitale et en cybersécurité, Wavestone constitue un terrain d'observation privilégié pour analyser la dualité de l'intégration de l'IA. En tant que consultant au sein de la practice Cybersécurité, confronté quotidiennement à l'intensité des missions de gestion de crise et d'évaluation de la maturité des organisations, l'opportunité nous est donnée d'observer directement la manière dont les équipes s'approprient — ou rejettent — ces nouveaux outils transformationnels.

\section{Énoncé de la problématique}
Face à ce constat, l'objectif de ce mémoire professionnel n'est pas d'analyser le processus technique de développement des outils d'IA, mais bien d'en étudier les impacts organisationnels et managériaux. Il s'agit de comprendre comment l'organisation navigue entre l'impératif stratégique d'innovation et la préservation de son capital humain et de sa fiabilité. 

Dès lors, la réflexion sera guidée par la problématique suivante : \textbf{L'intégration des outils d'intelligence artificielle générative dans le conseil : comment l'automatisation des tâches d'analyse redéfinit-elle la légitimité et le développement des compétences des consultants ?}